\documentclass[11pt]{amsart}
\usepackage{amsmath,amssymb,amsthm}
\usepackage[margin=1in]{geometry}

\newtheorem{theorem}{Theorem}[section]
\newtheorem{proposition}[theorem]{Proposition}
\newtheorem{lemma}[theorem]{Lemma}
\newtheorem{corollary}[theorem]{Corollary}
\theoremstyle{definition}
\newtheorem{definition}[theorem]{Definition}
\theoremstyle{remark}
\newtheorem{remark}[theorem]{Remark}

\newcommand{\ve}{\varepsilon}
\newcommand{\R}{\mathbb{R}}

\title{A quantitative version of the approximation theorem:\\ rate $O(\ve)$ for the high activation energy limit (working notes)}
\date{\today}

\begin{document}
\maketitle

\noindent\textbf{Summary.} We make Theorems 4.1--4.2 of the 2004 SIAM paper
(approximation of strict classical super/subsolutions of the free boundary
problem by super/subsolutions of $(P_\ve)$) quantitative, replacing the
blow-up argument by an explicit barrier: the inner profile composed with the
super/subsolution itself. The approximation error is $O(\ve)$ with a constant
that does \emph{not} depend on the strictness $\delta$, provided
$\ve\lesssim\delta$. Combined with the rescaling $v_\lambda$ of the 2004 paper
this gives
\[
\sup_{0\le t\le T}\|u^\ve(\cdot,t)-v(\cdot,t)\|_{L^\infty(\R^N)}\le C(\ve+\eta_\ve)
\]
whenever $(P)$ has a classical solution $v$ which is $C^{2,1}$ up to its free
boundary with non-degenerate gradient, in the geometry of the 2004 paper
(strictly starshaped data); here $\eta_\ve=\|w^\ve/\ve-w_0\|_\infty$. This is
the rate observed numerically.

\noindent\textbf{Status.} The structure of the argument is complete. Three
ingredients are standard and only sketched: smooth dependence of the inner
profiles on parameters (Lemma~\ref{lem:inner}(iv)), the comparison principle
for $(P_\ve)$ with $C^{1,1}$-type super/subsolutions and min/max of them, and
the construction of the extension $\tilde v$.

% ---------------------------------------------------------------------------
\section{Setting}

$Q_T=\R^N\times(0,T)$. For $\ve>0$,
\begin{equation}\label{eq:Pe}\tag{$P_\ve$}
u_t=\Delta u-(u+w^\ve)f_\ve(u),\qquad f_\ve(s)=\ve^{-2}f(s/\ve),
\end{equation}
where $w^\ve$ is caloric, and the limit problem is
\begin{equation}\label{eq:P}\tag{$P$}
u_t=\Delta u\ \text{ in }\{u>0\},\qquad |\nabla u|=\sqrt{2M(w_0)}\ \text{ on }\partial\{u>0\},
\end{equation}
with
\[
B(\tau;W)=\int_{-W}^{\tau}(s+W)f(s)\,ds,\qquad M(W)=B(1;W).
\]
\textbf{Assumptions.}
\begin{itemize}
\item[(F)] $f$ is Lipschitz (constant $L_f$), nonincreasing, $f>0$ on $(-\infty,1)$,
$f=0$ on $[1,\infty)$.
\item[(W)] $w_0$ is caloric, $w_0\ge-1+\delta_1$, and $W^\ve:=w^\ve/\ve$ satisfies
$\|W^\ve-w_0\|_\infty\le\eta_\ve$ and $\|W^\ve\|_{C^{2,1}(\overline{Q_T})}\le C_W$.
\end{itemize}
We write $p=(G,W)$ for the parameters of the inner profiles, and $W(x,t)=W^\ve(x,t)$.

\textbf{Comparison.} For fixed $\ve$ the nonlinearity $(u+w^\ve)f_\ve(u)$ is
Lipschitz in $u$, so bounded weak super/subsolutions of \eqref{eq:Pe} (in
particular functions in $C^1_x$ with locally bounded $D^2_x$, $\partial_t$) are
ordered if they are ordered at $t=0$; the minimum of two supersolutions is a
supersolution and the maximum of two subsolutions is a subsolution.

% ---------------------------------------------------------------------------
\section{Inner profiles with parameters}

Fix $\gamma>0$ and a compact set $\mathcal P_\gamma$ of parameters $p=(G,W)$ with
\[
0<g_1\le G\le g_2,\qquad -1+\delta_1/2\le W\le \bar W,\qquad G^2\le 2M(W)-2\gamma
\]
(the \emph{supersolution regime}); and $\mathcal P^\gamma$ with the last
condition replaced by $G^2\ge2M(W)+2\gamma$ (the \emph{subsolution regime}).

\begin{lemma}[Inner profiles]\label{lem:inner}
\textbf{(a) Supersolution regime.} There are $\kappa_0,C,\lambda_0>0$ depending on
$f,\gamma,g_{1,2},\delta_1,\bar W$ such that for $p\in\mathcal P_\gamma$ and
$0\le\kappa\le\kappa_0$ there are $\theta=\theta(p,\kappa)\in(0,1+W)$ and
$\Psi(\cdot;p,\kappa)\in C^{1,1}(\R)$ with
\begin{enumerate}
\item[(i)] $\Psi(s)=s$ for $s\ge1-\theta$; on $(-\infty,1-\theta)$,
$G^2\Psi''+\kappa\Psi'=(\Psi+W)f(\Psi)$, and $\Psi(s)\to-W$ as $s\to-\infty$;
\item[(ii)] $0<\Psi'\le1$, $-W<\Psi\le s^++1$, and $|\Psi(s)-s^+|\le C$;
\item[(iii)] $|\Psi''|\le C\Psi'$ on $(-\infty,1-\theta)$; for every $S>0$,
$\Psi'\ge c(S)>0$ on $[-S,1-\theta]$; and $0<\Psi+W\le Ce^{\lambda_0s}$,
$\Psi'\le Ce^{\lambda_0 s}$ for $s\le0$;
\item[(iv)] $\Psi$ is $C^2$ in $p$ for $s<1-\theta$, $\Psi_p$ is continuous across
$s=1-\theta$ (and vanishes for $s\ge1-\theta$), and for $s\le0$
\[
|\Psi_G|+|\Psi_W+1|+|\Psi_{sp}|+|\Psi_{pp}|\le C(1+|s|)^2e^{\lambda_0 s},
\]
while all these quantities are bounded by $C$ on $[0,1-\theta]$.
\end{enumerate}
\textbf{(b) Subsolution regime.} For $p\in\mathcal P^\gamma$, $0\le\kappa\le\kappa_0$ and
$\bar q\in(0,1]$ there are $s_*=s_*(p,\kappa,\bar q)>-C$ and
$\Psi\in C^{1,1}((s_*,\infty))$ with $\Psi(s)=s$ for $s\ge1$,
$G^2\Psi''-\kappa\Psi'=(\Psi+W)f(\Psi)$ on $(s_*,1)$, $\Psi(s_*)=-W-\bar q$,
$\Psi'\ge c>0$ on $(s_*,\infty)$, $|\Psi''|\le C$, $|\Psi(s)-s^+|\le C$, and the
analogue of (iv) with bounds $C$.
\end{lemma}

\begin{proof}
(a) Phase plane. Write $P=\Psi'$ as a function of $\Psi$ along the orbit; then
$\frac{d}{d\Psi}\big(\tfrac{G^2}{2}P^2\big)=(\Psi+W)f(\Psi)-\kappa P$. The orbit leaving
the saddle $(-W,0)$ (eigenvalues $\frac{-\kappa\pm\sqrt{\kappa^2+4G^2f(-W)}}{2G^2}$; note
$f(-W)\ge f(1-\delta_1/2)>0$) satisfies
\begin{equation}\label{eq:en}
\tfrac{G^2}{2}P(\Psi)^2=B(\Psi;W)-\kappa\int_{-W}^{\Psi}P\,d\Psi .
\end{equation}
As in the travelling wave lemma, $P>0$ as long as $\Psi<1$ and
$P\le\sqrt{2B(\Psi;W)}/G$, so $0\le\kappa\int P\le\kappa C$. We look for
$\Psi_m=1-\theta$ with $P(\Psi_m)=1$, i.e.
$B(\Psi_m;W)=\frac{G^2}2+\kappa\int_{-W}^{\Psi_m}P$. The right-hand side lies in
$[\frac{G^2}2,\frac{G^2}2+C\kappa]\subset[\frac{G^2}{2},M(W)-\frac\gamma2]$ if $\kappa\le\gamma/(2C)$;
since $B(\cdot;W)$ is continuous, strictly increasing on $[-W,1]$ from $0$ to
$M(W)$, and the right-hand side depends continuously on $\Psi_m$, a solution
$\Psi_m\in(-W,1)$ exists, with $1-\Psi_m\ge c(\gamma)>0$ (because
$B(\Psi_m)\le M-\gamma/2$ and $B'(\cdot)\le C$). Translate the orbit so that
$\Psi(1-\theta)=1-\theta$ and set $\Psi(s)=s$ for $s\ge1-\theta$: since
$P(\Psi_m)=1$, $\Psi\in C^1$, and $\Psi''$ is bounded with a jump at $1-\theta$.
(ii): $0<P\le1$ on the orbit up to $\Psi_m$ because $P$ is increasing in
$\Psi$ there: $\frac{d}{d\Psi}(\frac{G^2}2P^2)=(\Psi+W)f(\Psi)-\kappa P>0$, since near the
saddle $(\Psi+W)f\simeq f(-W)(\Psi+W)$ and $\kappa P\simeq\kappa\lambda_+(\Psi+W)$ with
$\kappa\lambda_+<f(-W)$ for $\kappa$ small, while away from it
$(\Psi+W)f(\Psi)\ge c(\gamma)>0$ for $\Psi\le\Psi_m\le1-c(\gamma)$; the bound on $|\Psi-s^+|$
follows from the exponential tail and $P\ge c$ on $[\Psi(-S),\Psi_m]$.
(iii): near the saddle, $\Psi''=\frac1{G^2}\big((\Psi+W)f(\Psi)-\kappa P\big)$ and
$P\simeq\lambda_+(\Psi+W)$, so $|\Psi''|\le C(\Psi+W)\le CP$; away from it $P\ge c$
and $\Psi''$ is bounded. The tail estimates are the same as in the travelling
wave lemma (uniform decay rate $\lambda_0=\frac12\sqrt{f(1-\delta_1/2)}/g_2$).
(iv): $C^2$ dependence on $(G,W,\kappa)$ of the unstable manifold and of the
transversal intersection defining $\Psi_m$ (implicit function theorem,
$\partial_{\Psi_m}[B-\frac{G^2}{2}-\kappa\int P]=(\Psi_m+W)f(\Psi_m)-\kappa\ge c>0$).
In the tail, $\Psi+W=a(p)e^{\lambda(p)s}(1+o(1))$, whose $p$-derivatives carry
the factors $s$ and $s^2$. Continuity of $\Psi_p$ at $1-\theta$:
$\Psi_p(1-\theta^-)=(1-\Psi'(1-\theta))\,\partial_p\theta=0$.

(b) Now $\frac{G^2}{2}P^2=B(\Psi;W)+\big(\frac{G^2}2-M(W)\big)+\kappa\int_\Psi^1P$ for
$\Psi\le1$ (integrating backwards from $P(1)=1$), and the second term is
$\ge\gamma$. Hence $P\ge\sqrt{2\gamma}/g_2$ on the whole backward orbit, which
crosses $\Psi=-W-\bar q$ at a finite $s_*$, with $|s_*|\le C$.
\end{proof}

% ---------------------------------------------------------------------------
\section{Quantitative approximation of strict supersolutions}

\begin{definition}
A function $\tilde u\ge0$ on $\overline{Q_T}$ is a \emph{regular strict
supersolution with margins $(\delta,\delta_0,r,K_0)$} if, with
$\Omega=\{\tilde u>0\}$ and $\Gamma=\partial\Omega\cap Q_T$:
\begin{enumerate}
\item $\Delta\tilde u-\tilde u_t\le0$ in $\Omega$;
\item $|\nabla\tilde u|\le\sqrt{2M(w_0)}-\delta$ on $\Gamma$;
\item there is a function $\tilde v$ on $\overline{Q_T}$ with $\tilde v=\tilde u$ on
$\overline\Omega$, $\tilde v<0$ outside $\overline\Omega$, $\tilde v\le-\delta_0r$ at
parabolic distance $\ge r$ from $\Gamma$ in the complement of $\Omega$,
$\tilde u\ge\delta_0 r$ at distance $\ge r$ from $\Gamma$ in $\Omega$, and, in the
$r$-neighbourhood $N_r$ of $\Gamma$: $\tilde v\in C^{2,1}$, $|\nabla\tilde v|\ge\delta_0$,
$\|\tilde v\|_{C^{2,1}(N_r)}\le K_0$; elsewhere $\|\tilde v\|_{C^{2,1}}\le K_0$.
\end{enumerate}
(The extension $\tilde v$ exists, e.g., by reflection across $\Gamma$ in normal
coordinates, as soon as $\tilde u$ is $C^{2,1}$ up to $\Gamma$ with
$|\nabla\tilde u|\ge2\delta_0$ there.)
\end{definition}

\begin{proposition}[Quantitative Theorem 4.1]\label{prop:super}
Assume (F), (W). There are $c_*,C$ depending only on $f,\delta_1,C_W,\delta_0,r,K_0$
(\emph{not} on $\delta$) such that if $\tilde u$ is a regular strict
supersolution with margins $(\delta,\delta_0,r,K_0)$ and
\[
\ve+\eta_\ve\le c_*\delta,
\]
then
\[
V^\ve(x,t)=\ve\,\Psi\Big(\frac{\tilde v(x,t)}\ve;\,p(x,t),K\ve\Big)+\ve^2\tilde q,\qquad
p=(G,W),\quad G=\min\{|\nabla\tilde v|,G_\gamma(W)\},
\]
with $G_\gamma(W)^2=2M(W)-2\gamma$, $\gamma=c\delta$, and suitable constants
$K,\tilde q$, is a supersolution of \eqref{eq:Pe} in $Q_T$ with
\[
\|V^\ve-\tilde u\|_{L^\infty(Q_T)}\le C\ve .
\]
\end{proposition}

\begin{proof}
\emph{Parameters are admissible.} On $\Gamma$, $|\nabla\tilde v|^2\le2M(w_0)-2\delta\sqrt{2M}+\delta^2$.
In the region where the profile is not the identity, $\tilde v<\ve(1-\theta)\le\ve$,
and either we are within distance $C\ve/\delta_0$ of $\Gamma$, where
$|\nabla\tilde v|$ changes by $O(K_0\ve/\delta_0)$ and $|W-w_0|\le\eta_\ve+C\ve$, so
$G^2\le2M(W)-2\gamma$ with $\gamma=c\delta$ if $\ve+\eta_\ve\le c_*\delta$; or we are in
the burnt region at distance $\ge C\ve/\delta_0$, where $G$ is the clipped value.
In all cases $p\in\mathcal P_\gamma$.

\emph{Derivatives.} Write $s=\tilde v/\ve$ and $\Psi=\Psi(s;p,\kappa)$, $\kappa=K\ve$.
Then
\[
\nabla V=\Psi_s\nabla\tilde v+\ve\Psi_p\nabla p,\qquad
\Delta V=\Psi_s\Delta\tilde v+\frac{1}{\ve}\Psi_{ss}|\nabla\tilde v|^2+2\Psi_{sp}\nabla\tilde v\cdot\nabla p
+\ve\Psi_p\Delta p+\ve\Psi_{pp}\nabla p\cdot\nabla p,
\]
\[
V_t=\Psi_s\tilde v_t+\ve\Psi_p p_t .
\]
$V$ is $C^1$ in $x$ (by Lemma~\ref{lem:inner}(iv) $\Psi_s$ and $\Psi_p$ are continuous)
and its second derivatives are bounded, so it suffices to check the
inequality $\mathcal R:=\Delta V-V_t-(V+w^\ve)f_\ve(V)\le0$ a.e.

\emph{Reaction.} With $V=\ve(\Psi+\ve\tilde q)$,
$(V+w^\ve)f_\ve(V)=\ve^{-1}(\Psi+\ve\tilde q+W)f(\Psi+\ve\tilde q)$, and since $f$ is
nonincreasing and Lipschitz,
\begin{equation}\label{eq:react}
(\Psi+\ve\tilde q+W)f(\Psi+\ve\tilde q)\ \ge\ (\Psi+W)f(\Psi)-L_f\ve\tilde q(\Psi+W)+\ve\tilde qf(\Psi+\ve\tilde q).
\end{equation}

\emph{Region $s\ge1-\theta$.} Here $V=\tilde u+\ve^2\tilde q$ and $\Psi_p=0$. If
$s\ge1$, $f_\ve(V)=0$ and $\mathcal R=\Delta\tilde u-\tilde u_t\le0$. If
$1-\theta\le s<1$, $V+w^\ve=\ve(s+\ve\tilde q+W)>0$, so the reaction term is
nonnegative and again $\mathcal R\le\Delta\tilde u-\tilde u_t\le0$.

\emph{Region $s<1-\theta$.} Collecting terms, using
$G^2\Psi_{ss}+\kappa\Psi_s=(\Psi+W)f(\Psi)$ (where $G=|\nabla\tilde v|$; where $G$ is
clipped, see below), \eqref{eq:react} and $\Delta W-W_t=0$ (so that the term
$\ve\Psi_W(\Delta W-W_t)$ vanishes):
\[
\mathcal R\le\Psi_s(\Delta\tilde v-\tilde v_t)-K\Psi_s+T_p+L_f\tilde q(\Psi+W)-\tilde qf(\Psi+\ve\tilde q),
\]
\[
T_p=2\Psi_{sp}\nabla\tilde v\cdot\nabla p+\ve\big(\Psi_G(\Delta G-G_t)+\Psi_{pp}\nabla p\cdot\nabla p\big).
\]
In $\Omega$, $\Delta\tilde v-\tilde v_t\le0$; outside, $|\Delta\tilde v-\tilde v_t|\le K_0$.
By (W) and the bounds on $\tilde v$, $|\nabla p|+|\Delta G-G_t|\le C$ in $N_r$.

\emph{Layer part} $-S\le s<1-\theta$ ($S$ fixed below). By
Lemma~\ref{lem:inner}(iii),(iv), $|T_p|\le C\le C\Psi_s/c(S)$ and
$L_f\tilde q(\Psi+W)\le C\tilde q\le C\tilde q\Psi_s/c(S)$; hence $\mathcal R\le0$ if
$K\ge K_0+C(1+\tilde q)/c(S)$.

\emph{Tail} $s<-S$. Here $\Psi+W\le Ce^{\lambda_0 s}$, so $\Psi+\ve\tilde q\le-W+Ce^{-\lambda_0S}+\ve\tilde q$
and $f(\Psi+\ve\tilde q)\ge a:=\frac12f(1-\delta_1/2)$ if $S$ is large and $\ve$ small.
Also $|T_p|\le C(1+|s|)^2e^{\lambda_0s}\le C(1+S)^2e^{-\lambda_0S}$ and
$L_f\tilde q(\Psi+W)\le C\tilde qe^{-\lambda_0S}$. Fix $S$ so that
$C(1+S)^2e^{-\lambda_0S}\le a/4$ and $Ce^{-\lambda_0S}L_f\le a/4$; then with $\tilde q=1$,
$\mathcal R\le K_0\Psi_s-K\Psi_s+a/4+a/4-a\le0$.

\emph{Clipped parameters.} If $G\ne|\nabla\tilde v|$, then $s\le-c\delta_0 r/\ve$ (far
burnt region, or the part of $N_r$ at distance $\ge C\ve/\delta_0$ from $\Gamma$), and the
extra term $\frac1\ve(|\nabla\tilde v|^2-G^2)\Psi_{ss}$ is bounded by
$\frac C\ve e^{-c/\ve}\le a/4$ for $\ve$ small; it is absorbed in the tail estimate.

\emph{Error.} $|V-\tilde u|\le\ve\sup|\Psi(s)-s^+|+\ve^2\tilde q\le C\ve$, using
$\tilde u=\ve s^+$ where $\tilde v=\ve s$.
\end{proof}

\begin{proposition}[Quantitative Theorem 4.2]\label{prop:sub}
Assume (F), (W), and let $\tilde u$ be a regular strict \emph{subsolution} with
margins $(\delta,\delta_0,r,K_0)$ (inequalities (1), (2) reversed:
$\Delta\tilde u-\tilde u_t\ge0$ in $\Omega$, $|\nabla\tilde u|\ge\sqrt{2M}+\delta$ on
$\Gamma$). If $\ve+\eta_\ve\le c_*\delta$, then
\[
\underline V^\ve=\max\Big\{\ve\Psi\Big(\frac{\tilde v}\ve;p,K\ve\Big),\ -w^\ve-\ve^2\tilde q\Big\}
\]
(with the profiles of Lemma~\ref{lem:inner}(b), $\bar q=\ve\tilde q$, and the first
function extended by $-\infty$ for $s<s_*$) is a subsolution of \eqref{eq:Pe} with
$\|\underline V^\ve-\tilde u\|_\infty\le C\ve$.
\end{proposition}

\begin{proof}
$-w^\ve-\ve^2\tilde q$ is a subsolution: it is caloric and
$(V+w^\ve)f_\ve(V)=-\ve^2\tilde qf_\ve(V)\le0$. The first function, $V_1=\ve\Psi(\tilde v/\ve)$, is a subsolution for
$s\ge s_*$. For $s\ge1$, $V_1=\tilde u\ge\ve$, there is no reaction and
$\mathcal R=\Delta\tilde u-\tilde u_t\ge0$. For $s_*\le s<1$ the reaction is
$\ve^{-1}(\Psi+W)f(\Psi)$ exactly, and the computation of
Proposition~\ref{prop:super} with reversed signs gives
$\mathcal R\ge\Psi_s(\Delta\tilde v-\tilde v_t)+K\Psi_s-|T_p|$, where
$\Delta\tilde v-\tilde v_t\ge0$ in $\Omega$ and $\ge-K_0$ outside. Now
$\Psi_s\ge c>0$ on the whole range $s\ge s_*$, which is bounded below (there is
no tail), and $|T_p|\le C$; choose $K\ge K_0+C/c$. Since $\Psi(s_*)=-W-\ve\tilde q$,
$V_1=-w^\ve-\ve^2\tilde q$ exactly at $s=s_*$ and $V_1>-w^\ve-\ve^2\tilde q$ for
$s>s_*$; thus $\underline V^\ve=V_1$ for $s\ge s_*$, $=-w^\ve-\ve^2\tilde q$ otherwise,
it is continuous, and it is a subsolution (maximum of two subsolutions, each
used where it is the larger). Finally $|\underline V^\ve-\tilde u|\le C\ve$ by
Lemma~\ref{lem:inner}(b).
\end{proof}

% ---------------------------------------------------------------------------
\section{Rate of convergence to a classical solution}

\begin{theorem}\label{thm:rate}
Assume (F), (W), and the setting of the 2004 SIAM paper: $u_0$ Lipschitz with
compact support, $w_0$ satisfies the monotonicity (2.2) there, and in addition
\begin{itemize}
\item[(S$'$)] $\Omega_0=\{u_0>0\}$ is strictly starshaped with respect to $0$
($x\cdot\nu\ge c_0>0$ on $\partial\Omega_0$), $u_0(\lambda x)\ge u_0(x)$ for
$0<\lambda<1$, and $u_0$ decreases radially at a definite rate near the
boundary: $x\cdot\nabla u_0\le-c_0$ in $\Omega_0\cap\{\operatorname{dist}(x,\partial\Omega_0)<\rho_0\}$
(in particular $u_0\ge c\operatorname{dist}(x,\partial\Omega_0)$ there).
\end{itemize}
Let $v$ be a classical solution of \eqref{eq:P} on $[0,T^*]$ with data $u_0$,
which is $C^{2,1}$ up to its free boundary $\Gamma_v$ (one sided) with
$|\nabla v|\ge2\delta_0>0$ near $\Gamma_v$. Let $u^\ve$ solve \eqref{eq:Pe} with
initial data satisfying $|u_0^\ve-u_0|\le C_0\ve$ on $\Omega_0$ and
$u_0^\ve=-w^\ve(\cdot,0)$ outside $\Omega_0$ (up to $C_0\ve^2$). Then for every
$T<T^*$ there is $C$ such that
\[
\sup_{0\le t\le T}\|u^\ve(\cdot,t)-v(\cdot,t)\|_{L^\infty(\R^N)}\le C(\ve+\eta_\ve).
\]
\end{theorem}

\begin{proof}
Let $\ve'=\ve+\eta_\ve$ and, for a constant $a$ to be chosen, $\lambda=1-a\ve'$,
$\lambda'=1-a\ve'/2$. As in Section 2 of the 2004 paper,
$v_\lambda(x,t)=\frac1{\lambda'}v(\lambda x,\lambda^2t)$ is a classical supersolution on
$[0,T^*/\lambda^2]\supset[0,T]$ with
$|\nabla v_\lambda|\le\sqrt{2M(w_0)}-\delta$ on its free boundary,
$\delta=(1-\lambda/\lambda')\sqrt{2M_0}\ge c\,a\ve'$ (this uses (2.2), i.e.
$M(\lambda x,\lambda^2t)\le M(x,t)$). It inherits the regularity of $v$: it is a
regular strict supersolution with margins $(\delta,\delta_0,r,K_0)$, with
$\delta_0,r,K_0$ independent of $\ve$. Choose $a\ge1/(cc_*)$, so that
$\ve+\eta_\ve\le c_*\delta$, and apply Proposition~\ref{prop:super}: $V^\ve$ is a
supersolution with $|V^\ve-v_\lambda|\le C\ve$.

\emph{Initial ordering.} We claim $v_\lambda(x,0)-u_0(x)\ge c\,a\ve'$ on
$\overline\Omega_0$. Near $\partial\Omega_0$ (distance $<\rho_0$), by (S$'$),
$u_0(\lambda x)-u_0(x)=-\int_\lambda^1x\cdot\nabla u_0(\tau x)\,d\tau\ge c_0(1-\lambda)$ (the
segment $\{\tau x\}$ stays in the region where $x\cdot\nabla u_0\le-c_0$, or enters the
interior where $u_0$ is larger); away from $\partial\Omega_0$, $u_0\ge c$ and
$v_\lambda(x,0)-u_0(x)\ge(\frac1{\lambda'}-1)u_0(\lambda x)\ge c\,a\ve'$. Hence, on $\Omega_0$,
$V^\ve(x,0)\ge v_\lambda(x,0)-C\ve\ge u_0(x)+(ca-C)\ve'\ge u_0^\ve(x)$ if $a$ is large. Outside
$\Omega_0$, $V^\ve\ge-w^\ve+\ve^2\tilde q\ge u_0^\ve$ (by Lemma~\ref{lem:inner}(ii),
$\Psi>-W$). By comparison, $u^\ve\le V^\ve\le v_\lambda+C\ve$, and
$v_\lambda-v\le C(1-\lambda)\le Ca\ve'$ on $[0,T]$ ($v$ is Lipschitz in $x$ and
$1/2$-H\"older in $t$, and bounded). Thus $u^\ve-v\le C\ve'$.

\emph{Lower bound.} Same argument with $\lambda=1+a\ve'$, $\lambda'=1+a\ve'/2$:
$v_\lambda$ is a regular strict subsolution with margin $\delta\ge ca\ve'$
(now $M(\lambda x,\lambda^2t)\ge M(x,t)$) and, by the same computation,
$u_0-v_\lambda(\cdot,0)\ge c\,a\ve'$ on $\overline{\lambda^{-1}\Omega_0}$, while on
$\Omega_0\setminus\lambda^{-1}\Omega_0$ (a strip of width $\ge c(1-1/\lambda)$ along
$\partial\Omega_0$) $v_\lambda(\cdot,0)=0$ and $u_0\ge0$. Proposition~\ref{prop:sub} gives a subsolution
$\underline V^\ve\le u^\ve$ (initial ordering as above, using
$\underline V^\ve=-w^\ve-\ve^2\tilde q\le u_0^\ve$ in the burnt region), and
$u^\ve\ge\underline V^\ve\ge v_\lambda-C\ve\ge v-C\ve'$.
\end{proof}

\begin{remark}
\begin{enumerate}
\item The only place where $\delta$ enters is the admissibility of the
parameters, $\ve+\eta_\ve\le c_*\delta$; the constant of Proposition~\ref{prop:super}
does not depend on $\delta$. This is what allows $1-\lambda\sim\ve$ in the
rescaling argument, and gives the rate $\ve$ rather than $\sqrt\ve$.
\item The theorem gives a rate for the convergence proved (by compactness) in
the 2000 IFB paper and for the uniqueness of Theorem 5.3 of the 2004 paper,
under the additional regularity of $v$. The rate $O(\ve)$ is sharp: the
level set $\{u^\ve=\ve\}$ is at distance $\ve/\sqrt{2M}+o(\ve)$ from $\Gamma_v$
(travelling waves, and the experiments).
\item Combined with the discrete estimates of the convergence notes, this is
the missing link for general data: the discrete solution approximates $u^\ve$
(away from extinction and with a front that the mesh follows), and $u^\ve$
approximates $v$ at rate $\ve$.
\end{enumerate}
\end{remark}

\end{document}
